
\begin{obeassessment}
\textbf{Kuis Bab 4}
1. Mengapa Vigen\`ere Cipher lebih aman daripada Caesar Cipher terhadap serangan analisis frekuensi?
2. Jelaskan mengapa dalam Playfair Cipher huruf 'J' biasanya dihilangkan atau digabung dengan 'I'.
3. Apa syarat utama agar sebuah matriks dapat digunakan sebagai kunci dalam Hill Cipher?
4. Jelaskan mengapa parameter $m$ pada affine cipher harus relatif prima dengan ukuran alfabet, dan berikan satu contoh nilai $m$ yang tidak sah beserta akibatnya bagi proses dekripsi.
5. Diberikan cipherteks Vigen\`ere dengan kriptogram berulang berjarak 15, 20, dan 30 huruf. Tentukan panjang kunci yang paling mungkin dan uraikan langkah-langkahnya.
6. Sebutkan tiga syarat kerahasiaan sempurna One-Time Pad, lalu jelaskan secara singkat apa yang terjadi bila masing-masing syarat dilanggar.
\end{obeassessment}
\begin{competencychecklist}
\begin{tabularx}{\linewidth}{|X|c|}
\hline
Indikator Kompetensi & Tercapai \\
\hline
Saya dapat membedakan cipher substitusi dan transposisi & [ ] \\
\hline
Saya dapat melakukan analisis frekuensi untuk memecahkan monoalphabetic cipher & [ ] \\
\hline
Saya dapat melakukan enkripsi dan dekripsi Vigen\`ere Cipher & [ ] \\
\hline
Saya dapat menerapkan aturan enkripsi Playfair Cipher & [ ] \\
\hline
Saya dapat menghitung operasi matriks pada Hill Cipher & [ ] \\
\hline
Saya dapat menjelaskan prinsip kerja rotor pada mesin Enigma & [ ] \\
\hline
Saya dapat melakukan enkripsi dan dekripsi Affine Cipher serta menentukan ruang kuncinya & [ ] \\
\hline
Saya dapat melakukan kriptanalisis Hill dan Affine Cipher dengan serangan \textit{known-plaintext} & [ ] \\
\hline
Saya dapat menentukan panjang kunci Vigen\`ere dengan metode Kasiski & [ ] \\
\hline
Saya dapat menjelaskan kerahasiaan sempurna One-Time Pad beserta tiga syaratnya & [ ] \\
\hline
Saya dapat menjelaskan mengapa One-Time Pad tidak sesuai untuk komunikasi modern & [ ] \\
\hline
\end{tabularx}
\end{competencychecklist}
